% Формат файла
\documentclass[12pt]{article}
\usepackage[paperheight=297mm,
	paperwidth=210mm,
	top=20mm,
	bottom=20mm,
	left=15mm,
right=15mm]{geometry}

% Импорт преамбулы
\usepackage{../mypreamble}

% Оглавление
\title{\textbf{Алгоритмы и структуры данных}}\date{}\author{}

\begin{document}

\maketitle
\tableofcontents
\label{toc}
\newpage

\section{Сложность алгоритмов}

\begin{definition}[Оценка «O-большое»]
	\[f(n)=O(g(n))\iff \exists\,N_0\colon \forall
	n>N_0\ \exists\,c\colon f(n)\leqslant c\cdot g(n)\]
\end{definition}
\begin{definition}[Оценка «O-малое]
	\[f(n)=o(g(n))\iff \exists\,N_0\colon \forall
	n>N_0\ \exists\,c\colon f(n)< c\cdot g(n)\]
\end{definition}
\begin{definition}[Оценка «Омега-большое»]
	\[f(n)=\Omega(g(n))\iff \exists\,N_0\colon \forall
	n>N_0\ \exists\,c\colon f(n)\geqslant c\cdot g(n)\]
\end{definition}
\begin{definition}[Оценка «Омега-малое»]
	\[f(n)=\omega(g(n))\iff \exists\,N_0\colon \forall
	n>N_0\ \exists\,c\colon f(n)> c\cdot g(n)\]
\end{definition}
\begin{definition}[Оценка «Тета»]
	\[f(n)=\Theta(g(n))\iff \exists\,N_0\colon \forall
		n>N_0\ \exists\,c_1, c_2\colon c_1\cdot g(n)\leqslant f(n)\leqslant
	c_2\cdot g(n)\]
\end{definition}
\end{document}